% Route values, deterministic choice, and conditional sorting.
\section{Route Values and Organizational Sorting}
\label{sec:p04-route-sorting}

This section composes the commercial-value kernel with the internal and
entrusted manufacturing technologies defined above. It treats the qualified
manufacturing-service price \(p_m\) as fixed and conjectured. Market
consistency for \(p_m^*\) is imposed in the capacity-market section below.

\subsection{Route values}
\label{subsec:p04-route-values}

The economic value of internal production is
\begin{equation}
  W_i^I(q,m)
  =
  s(q)R\!\left(q,c_I(m,k_i)\right)
  -
  F_I(m,k_i).
  \label{eq:p04-internal-value}
\end{equation}
Its financing-adjusted value is
\begin{equation}
  \widetilde W_i^I(q,m;\ell_i)
  =
  \begin{cases}
    W_i^I(q,m), & J_I(m,k_i)\leq\ell_i,\\
    -\infty, & J_I(m,k_i)>\ell_i.
  \end{cases}
  \label{eq:p04-finance-adjusted-internal}
\end{equation}

At a conjectured CMO capacity price \(p_m\), retained entrusted manufacturing
has value
\begin{equation}
  \begin{aligned}
  W_i^E(q,m;M,p_m)
  ={}&
  s(q)R\!\left(q,c_E(m)\right)
  -F_E(m)-p_m b(m)\\
  &-\mu_E-\tau_E(M).
  \end{aligned}
  \label{eq:p04-entrusted-value}
\end{equation}
The developer remains the authorization holder under \(E\); the qualified
external producer manufactures. Each route cost is subtracted exactly once.
In particular, \(p_m b(m)\) is not contained in \(c_E(m)\).
The financing-adjusted entrusted value is
\begin{equation}
  \widetilde W_i^E(q,m;M,p_m,\ell_i)
  =
  \begin{cases}
    W_i^E(q,m;M,p_m), & J_E(m)\leq\ell_i,\\
    -\infty, & J_E(m)>\ell_i.
  \end{cases}
  \label{eq:p04-finance-adjusted-entrusted}
\end{equation}
Neither \(J_I\) nor \(J_E\) is subtracted from route value: each is a
liquidity threshold, not a second copy of a real cost.

Let the noncore transfer/out-license outside value be a finite continuous
function \(T(q,m)\), and normalize abandonment or indefinite delay to zero:
\begin{equation}
  W^T(q,m)=T(q,m),
  \qquad
  W^A=0.
  \label{eq:p04-outside-values}
\end{equation}
No separate transfer market is added.

\subsection{Deterministic route choice}
\label{subsec:p04-deterministic-choice}

The optimized value of a planning-stage project is
\begin{equation}
  \widetilde W_i(q,m;M,p_m,\ell_i)
  =
  \max\left\{
    \widetilde W_i^I(q,m;\ell_i),
    \widetilde W_i^E(q,m;M,p_m,\ell_i),
    W^T(q,m),
    W^A
  \right\}.
  \label{eq:p04-optimized-value}
\end{equation}
The route choice is
\begin{equation}
  r_i^*(q,m;M,p_m,\ell_i)
  =
  \arg\max_{r\in\{I,E,T,A\}}
  \begin{cases}
    \widetilde W_i^I(q,m;\ell_i), & r=I,\\
    \widetilde W_i^E(q,m;M,p_m,\ell_i), & r=E,\\
    W^T(q,m), & r=T,\\
    W^A, & r=A.
  \end{cases}
  \label{eq:p04-route-choice}
\end{equation}
Continuous project and developer heterogeneity makes exact ties a
measure-zero event, so the choice is single valued almost surely. This is a
deterministic comparison; it generates no probabilistic route share.

When \(M=0\), \(\tau_E(0)=+\infty\) implies
\(\widetilde W_i^E=-\infty\).
At fixed \(p_m\), the binary reform comparison is therefore
\begin{equation}
  \begin{aligned}
  W_i^0(q,m;\ell_i)
  &=\max\{\widetilde W_i^I(q,m;\ell_i),W^T(q,m),0\},\\
  W_i^1(q,m;p_m,\ell_i)
  &=\max\{W_i^0(q,m;\ell_i),
    \widetilde W_i^E(q,m;1,p_m,\ell_i)\},\\
  W_i^1-W_i^0
  &=[\widetilde W_i^E-W_i^0]_+\geq0.
  \end{aligned}
  \label{eq:p04-binary-value-effect}
\end{equation}
The gain is exactly zero, and MAH has zero project-value and route-choice
effect, whenever the newly
feasible entrusted value fails to exceed the best pre-reform route.

\subsection{Internal--entrusted value gap}
\label{subsec:p04-value-gap}

For fixed \(q,m,M,p_m\), define
\begin{align}
  \Delta_{IE}(k_i;q,m,M,p_m)
  &=
  W_i^I(q,m)-W_i^E(q,m;M,p_m) \notag\\
  &=
  s(q)\!\left[
    R\!\left(q,c_I(m,k_i)\right)
    -R\!\left(q,c_E(m)\right)
  \right]
  -F_I(m,k_i) \notag\\
  &\quad
  +F_E(m)+p_m b(m)+\mu_E+\tau_E(M).
  \label{eq:p04-value-gap}
\end{align}
Conditional on both retained routes being financially feasible, only the
internal technology terms in the economic value gap vary with \(k_i\).
Holding \(q,m,M,p_m\) fixed,
\begin{equation}
  \frac{\partial\Delta_{IE}}{\partial k_i}
  =
  s(q)R_c\!\left(q,c_I(m,k_i)\right)c_{I,k}(m,k_i)
  -
  F_{I,k}(m,k_i)
  >0.
  \label{eq:p04-gap-slope}
\end{equation}
The first term is weakly positive because \(s(q)\geq0\),
\(R_c<0\), and \(c_{I,k}<0\); the second is strictly positive because
\(F_{I,k}<0\). Hence the entire derivative is strictly positive, including
the boundary case \(s(q)=0\).

\subsection{Conditional organizational cutoff}
\label{subsec:p04-cutoff}

Consider a finite entrusted barrier and fix \(\ell_i\) such that route \(E\)
is financeable. Since \(J_{I,k}<0\), the set of capabilities at which route
\(I\) is also financeable is an upper interval. On a connected portion of
that common-financeability set, suppose the continuous value gap satisfies the
endpoint crossing conditions
\[
  \Delta_{IE}(k_L)<0,
  \qquad
  \Delta_{IE}(k_H)>0,
  \qquad k_L<k_H.
\]
Continuity gives existence of a root, and
\eqref{eq:p04-gap-slope} gives strict monotonicity and therefore uniqueness.
Denote the unique root by
\begin{equation}
  \Delta_{IE}\!\left(k^*(q,m;p_m,M);q,m,M,p_m\right)=0.
  \label{eq:p04-cutoff}
\end{equation}
Conditional on both routes being financeable and both dominating transfer and
abandonment,
\begin{equation}
  k_i<k^*\Rightarrow r_i^*=E,
  \qquad
  k_i>k^*\Rightarrow r_i^*=I.
  \label{eq:p04-cutoff-sorting}
\end{equation}
If \(T\) or \(A\) dominates, the \(I/E\) cutoff does not determine the
realized route.

For a finite wedge, let the arguments \(\tau_E\) and \(p_m\) vary locally
while holding \(q,m\) fixed. Since
\[
  \frac{\partial\Delta_{IE}}{\partial\tau_E}=1,
  \qquad
  \frac{\partial\Delta_{IE}}{\partial p_m}=b(m)>0,
\]
the implicit-function theorem yields
\begin{equation}
  \frac{\partial k^*}{\partial\tau_E}
  =
  -\frac{1}{\Delta_{IE,k}(k^*)}<0,
  \qquad
  \frac{\partial k^*}{\partial p_m}
  =
  -\frac{b(m)}{\Delta_{IE,k}(k^*)}<0.
  \label{eq:p04-cutoff-derivatives}
\end{equation}
A lower finite institutional barrier raises \(k^*\), expanding the
entrusted-production region \(k_i<k^*\). A higher fixed CMO price lowers
\(k^*\), contracting that region. The second derivative is a fixed-price
comparative static, not an equilibrium derivative through \(p_m^*\).

\subsection{Financing corridor and MAH-relevant set}
\label{subsec:p04-financing-corridor}

The financing-adjusted MAH-relevant set is
\begin{equation}
  \mathcal C_i(p_m,\ell_i)
  =
  \left\{(q,m):
  J_E(m)\leq\ell_i,\quad
  W_i^E(q,m;1,p_m)>W_i^0(q,m;\ell_i)
  \right\}.
  \label{eq:p04-finance-relevant-set}
\end{equation}
The fixed-price gain is strict exactly on this set.

For a fixed \((q,m,k_i,p_m)\), suppose the locally relevant capital ordering
\begin{equation}
  J_E(m)<J_I(m,k_i)
  \label{eq:p04-capital-ordering}
\end{equation}
holds. This is a proposition-specific condition for research-oriented or
low-internal-capability developers, not a global assumption. Financing
capacity then creates three regions:
\begin{enumerate}
  \item \(\ell_i<J_E(m)\): neither retained route is financeable and MAH has
  no project-value effect;
  \item \(J_E(m)\leq\ell_i<J_I(m,k_i)\): entrusted production is financeable
  but internalization is not. The MAH gain is strict if and only if
  \(W_i^E>\max\{W^T,0\}\);
  \item \(\ell_i\geq J_I(m,k_i)\): both retained routes are financeable and
  the gain is \([W_i^E-\max\{W_i^I,W^T,0\}]_+\).
\end{enumerate}
The treatment effect can therefore rise at the lower corridor boundary and
fall, remain positive, or become zero when internalization returns to the old
choice set. It has no global monotonic sign in \(\ell_i\). By contrast, within
either fixed regime, a larger \(\ell_i\) weakly expands the feasible choice
set and weakly raises optimized project value.

\subsection{Organizational sorting result}
\label{subsec:p04-sorting-proposition}

Suppose the commercial-value and technology restrictions above hold, both
retained routes are financeable, the entrusted barrier is finite, the
endpoint crossing conditions hold on the common-financeability interval, and
ties have measure zero. Then the internal-minus-entrusted economic value gap
is strictly increasing in developer manufacturing capability, a unique finite
\(I/E\) cutoff exists, and the conditional sorting and implicit derivatives in
\eqref{eq:p04-cutoff-sorting}--\eqref{eq:p04-cutoff-derivatives} follow.
The proposition is conditional: it does not displace the valid transfer and
abandonment outside options.

\subsection{Zero-effect and non-tautology boundary}
\label{subsec:p04-zero-effect}

The route comparison is not tautological. Internal value and its financing
threshold depend on
developer capability through \(c_I\) and \(F_I\); entrusted value instead uses
qualified-external technology, a physical capacity requirement, a retained
holder burden, its own financing threshold, and the institutional wedge. MAH
can nevertheless have zero effect when \(E\) is not financeable, remains below
the pre-reform maximum, when \(T\) or \(A\) dominates, or when the relevant
project--firm pair lies outside the finite cutoff-crossing region. No smooth
random-utility term, inclusive-value object, borrowing decision, or
probabilistic route share is used.
